\chapter{为什么要建模？}

别担心。不会再有那种围绕“你抬头看会不会影响别人抬头看”的章节导言了。那个例子已经被我们榨到了几乎不可能再榨的地步。本书第二部分的例子大多使用真实数据，数据集可以从本书网站下载。

第一部分主要是概念性的，计算被压到最低限度，用手就能完成。相比之下，第二部分介绍的内容需要借助计算机来拟合回归模型，例如线性模型与 logistic 模型。本书无法详细介绍回归技术，因此我们假定读者对这类常用模型已有基本的理解和使用经验。本书网站提供了 R、SAS、Stata 与 Python 代码的链接，可复现正文所述的分析。代码边注标出代码中与正文分析相关的那一部分。

本章说明第一部分所用的非参数估计量与第二部分所用的参数（基于模型的）估计量有何差别，也回顾平滑的概念，并简要讨论任何建模决策都涉及的偏倚—方差权衡。本章要说明的是：无论分析目标是因果推断还是（比如说）预测，为什么数据分析都需要模型。在下一章之前，我们暂且搁置因果方面的考虑。请记住，关于建模的统计学文献浩如烟海，本章只能点出其中一些关键问题。

\section{数据不会自己说话}

设想一个由 16 名感染人类免疫缺陷病毒（HIV）的个体构成的研究总体。与本书第一部分不同，我们不再把这 16 人看作 10 亿个与他们完全相同者的代表。他们只是从一个很大、或许只是假想的超总体——目标总体——中随机抽取的 16 个人。

研究开始时，每个个体接受某一水平的治疗 $A$（抗逆转录病毒治疗），并在整个研究期间保持不变。研究结束时，测量所有个体的连续结局 $Y$（CD4 细胞计数，单位为 cells/mm$^3$）。我们希望一致地估计：在抽取这 16 人的总体中，治疗水平为 $A=a$ 者的 $Y$ 均值。也就是说，估计目标是未知的总体参数 $\E[Y\given A=a]$。

如第 10 章所定义，$\E[Y\given A=a]$ 的估计量 $\widehat{\E}[Y\given A=a]$ 是用来估计该未知总体参数的某个数据函数。粗略地说，一致估计量 $\widehat{\E}[Y\given A=a]$ 满足“样本量越大，估计值越接近总体值 $\E[Y\given A=a]$”这一要求。\mnote{一致估计量的严格定义见第 10 章。}估计量 $\widehat{\E}[Y\given A=a]$ 的两个例子是：(i) 接受 $A=a$ 者中 $Y$ 的样本均值；(ii) 数据集中恰好取 $A=a$ 的第一个观测的取值。接受 $A=a$ 者中 $Y$ 的样本均值是总体均值的一致估计量；$A=a$ 的第一个观测的取值则不是。实践中我们要求所有估计量都一致，因此用样本均值来估计总体均值。

假设治疗 $A$ 是二值变量，有两个可能取值：未治疗（$A=0$）与治疗（$A=1$）。一半个体接受了治疗（$A=1$）。图 11.1 是散点图，把 16 个个体各画成一个点，点的高度表示该个体结局 $Y$ 的取值。8 名受治疗者排在 $A=1$ 这一列，8 名未受治疗者排在 $A=0$ 这一列。如第 10 章所定义，总体中治疗水平为 $A=a$ 者的 $Y$ 均值之估计，就是把估计量——在本例中是样本均值——作用于某一具体数据集所得的数值结果。

\begin{figure}[htbp]
\centering
\begin{tikzpicture}[x=2.3cm,y=0.0152cm,causal]
  \draw[->] (0,0) -- (0,272);
  \draw[->] (0,0) -- (1.62,0);
  \foreach \y in {50,100,150,200,250}{%
    \draw (-0.045,\y) -- (0,\y);
    \node[left=2pt,font=\footnotesize,inner sep=1.5pt] at (0,\y) {\y};}
  \node[left=2pt,font=\footnotesize,inner sep=1.5pt] at (0,0) {0};
  \foreach \x/\lab in {0.55/0,1.35/1}{%
    \draw (\x,0) -- (\x,-8);
    \node[below=1pt,font=\small,inner sep=1pt] at (\x,-8) {\lab};}
  \node[below=13pt,font=\small] at (0.95,-8) {$A$};
  \foreach \yy in {170,110,80,65,50,30,22,15}{%
    \fill (0.55,\yy) circle[radius=1.7pt];}
  \foreach \yy in {220,200,180,168,152,110,90,50}{%
    \fill (1.35,\yy) circle[radius=1.7pt];}
\end{tikzpicture}
\caption{16 名个体的结局取值，治疗 $A$ 为二值（0：未治疗，1：治疗）。}
\label{fig:11.1}
\end{figure}

\mnote{图 11.1}我们对受治疗者总体均值的估计，是 $A=1$ 者的样本均值 146.25；对未受治疗者总体均值的估计，是 $A=0$ 者的样本均值 67.50。在受治疗者与未受治疗者可交换的前提下，差值 $146.25-67.50$ 可解释为目标总体中治疗 $A$ 对结局 $Y$ 的平均因果效应之估计。然而本章讲的不是做因果推断。我们当前的目标只是说明：在估计 $\E[Y\given A=a]$ 这类总体量时，为什么需要模型——无论估计值是否具有因果解释。

假设治疗 $A$ 是多值变量，可取 4 个值：未治疗（$A=1$）、低剂量治疗（$A=2$）、中剂量治疗（$A=3$）与高剂量治疗（$A=4$）。每个治疗水平各有四分之一的人接受。图 11.2 显示研究总体中 16 人的结局取值。为估计由治疗水平所定义的 4 组的总体均值，我们计算相应的样本均值；对 $A=1,2,3,4$，估计值依次为 70.0、80.0、117.5 与 195.0。\mnote{代码：程序 11.1}

\begin{figure}[htbp]
\centering
\begin{tikzpicture}[x=1.15cm,y=0.0152cm,causal]
  \draw[->] (0,0) -- (0,272);
  \draw[->] (0,0) -- (4.45,0);
  \foreach \y in {50,100,150,200,250}{%
    \draw (-0.09,\y) -- (0,\y);
    \node[left=2pt,font=\footnotesize,inner sep=1.5pt] at (0,\y) {\y};}
  \node[left=2pt,font=\footnotesize,inner sep=1.5pt] at (0,0) {0};
  \foreach \x in {0,1,2,3,4}{%
    \draw (\x,0) -- (\x,-8);
    \node[below=1pt,font=\small,inner sep=1pt] at (\x,-8) {\x};}
  \node[below=13pt,font=\small] at (2,-8) {$A$};
  \foreach \yy in {40,50,80,110}{%
    \fill (1,\yy) circle[radius=1.7pt];}
  \foreach \yy in {30,50,70,170}{%
    \fill (2,\yy) circle[radius=1.7pt];}
  \foreach \yy in {50,110,130,180}{%
    \fill (3,\yy) circle[radius=1.7pt];}
  \foreach \yy in {150,200,210,220}{%
    \fill (4,\yy) circle[radius=1.7pt];}
\end{tikzpicture}
\caption{16 名个体的结局取值，治疗 $A$ 有 4 个水平（1：未治疗，2：低剂量，3：中剂量，4：高剂量）。}
\label{fig:11.2}
\end{figure}

图 11.1 与图 11.2 分别是 2 类与 4 类离散（分类）变量的例子。由于研究个体数固定为 16，类别越多，每类的个体数就越少。各类别的样本均值仍是相应总体均值的无偏估计量（而且是精确无偏的），但各类个体数越少，样本均值接近相应总体均值的概率就越低。对图 11.2 的数据，各分类均值之 95\% 置信区间（见第 10 章）的长度将大于图 11.1 的数据。\mnote{图 11.2}

最后，设 $A$ 表示以 mg/日计的治疗剂量，取 0 到 100 mg 的整数值。图 11.3 显示 16 名个体各自的结局取值。由于治疗可能取值的个数远多于研究中的个体数，许多 $A$ 值无人接受。例如，研究总体中没有人接受治疗剂量 $A=90$。\mnote{图 11.3}

\begin{figure}[htbp]
\centering
\begin{tikzpicture}[x=0.062cm,y=0.0152cm,causal]
  \draw[->] (0,0) -- (0,272);
  \draw[->] (0,0) -- (104,0);
  \foreach \y in {50,100,150,200,250}{%
    \draw (-1.5,\y) -- (0,\y);
    \node[left=1pt,font=\footnotesize,inner sep=1.5pt] at (0,\y) {\y};}
  \node[left=1pt,font=\footnotesize,inner sep=1.5pt] at (0,0) {0};
  \foreach \x in {0,10,20,30,40,50,60,70,80,90,100}{%
    \draw (\x,0) -- (\x,-8);
    \node[below=1pt,font=\scriptsize,inner sep=0.8pt] at (\x,-8) {\x};}
  \node[below=13pt,font=\small] at (50,-8) {$A$};
  \foreach \xa/\ya in {2/22,11/55,13/10,18/33,24/102,30/85,38/62,44/157,%
                       44/188,53/120,59/228,66/112,70/215,79/200,84/140,97/218}{%
    \fill (\xa,\ya) circle[radius=1.7pt];}
\end{tikzpicture}
\caption{16 名个体的结局取值，治疗 $A$ 为 0 到 100 mg/日之间的剂量。}
\label{fig:11.3}
\end{figure}

这就产生一个问题：如何估计目标总体中治疗水平为 $A=90$ 者的结局 $Y$ 之均值？我们在图 11.1 与图 11.2 的数据中所用的估计量——按治疗水平分组的样本均值——在图 11.3 中无人取值的那些治疗水平上没有定义。若治疗 $A$ 真正连续，则几乎所有治疗水平上的样本均值都无定义。（连续变量 $A$ 可以看作有无穷多个类别的分类变量。）

以上描述表明，我们并不总能“让数据自己说话”而得到有意义的估计。相反，我们往往需要用一个模型来补充数据，下一节即讨论这一点。

\section{条件均值的参数估计量}

我们想由图 11.3 的数据估计治疗水平为 $A=90$ 者中 $Y$ 的均值，即 $\E[Y\given A=90]$。设我们预期 $A=90$ 者的 $Y$ 均值介于 $A=80$ 者的均值与 $A=100$ 者的均值之间。事实上，设想我们知道：在 $A$ 的整个取值范围内，$Y$ 的按治疗分组的总体均值都是治疗取值 $A$ 的线性函数。更确切地说，我们知道 $Y$ 的均值 $\E[Y\given A]$ 从 $A=0$ 时的某个值 $\theta_0$ 出发，每增加一个单位的 $A$ 就增加（或减少）$\theta_1$ 个单位。或者更紧凑地写成
\[
\E[Y\given A]=\theta_0+\theta_1 A
\]
这个方程是对条件均值函数 $\E[Y\given A]$ 形状的一种限制。\mnote{更一般地说，对关系形状的这种限制称为\emph{函数形式}（functional form），有些作者称之为\emph{剂量—反应曲线}；我们不用后一说法，因为它暗示治疗剂量会因果地影响反应，而在存在混杂时这可能是错的。}这一特定的限制称为\emph{线性均值模型}，量 $\theta_0$ 与 $\theta_1$ 称为模型的\emph{参数}。用有限个参数描述条件均值函数的模型，称为\emph{参数条件均值模型}。在本例中，参数 $\theta_0$ 与 $\theta_1$ 定义了一条直线，它在 $\theta_0$ 处与纵轴相交（截距），斜率为 $\theta_1$。也就是说，该模型设定所有条件均值函数都是直线，只是截距与斜率可以不同。

现在可以把图 11.3 的数据与参数均值模型结合起来，估计 $a$ 从 0 到 100 的所有 $\E[Y\given A=a]$。第一步是得到参数 $\theta_0$ 与 $\theta_1$ 的估计 $\hat\theta_0$ 与 $\hat\theta_1$；第二步是用这些估计去估计任一 $A=a$ 处 $Y$ 的均值。例如，要估计治疗水平 $A=90$ 者的 $Y$ 均值，我们使用表达式 $\widehat{\E}[Y\given A=90]=\hat\theta_0+90\hat\theta_1$。对每个个体的估计 $\widehat{\E}[Y\given A]$ 称为\emph{预测值}。

$\theta_0$ 与 $\theta_1$ 的精确无偏估计可由普通最小二乘法得到。下面给出该方法的非技术性说明。考虑图 11.3 的所有候选直线，每条对应截距 $\theta_0$ 与斜率 $\theta_1$ 的一个不同组合。对每条候选直线，可以计算每个点到该直线的垂直距离（残差），把这 16 个残差各自平方，再把 16 个平方残差相加。平方和最小的那条线就是“最小二乘”线，其参数值 $\hat\theta_0$ 与 $\hat\theta_1$ 就是“最小二乘”估计。$\hat\theta_0$ 与 $\hat\theta_1$ 可用线性代数轻易算出，任何统计学教科书都有介绍。\mnote{图 11.4}

\mnote{代码：程序 11.2。在残差方差不依赖于 $A$（同方差）的假设下，Wald 95\% 置信区间为：$\theta_0$ 的 $(-21.2,\,70.3)$，$\theta_1$ 的 $(1.28,\,2.99)$，以及 $\E[Y\given A=90]$ 的 $(172.1,\,261.6)$。}

在本例中，参数估计为 $\hat\theta_0=24.55$ 与 $\hat\theta_1=2.14$，它们定义了图 11.4 所示的直线。因此治疗水平 $A=90$ 者中 $Y$ 的预测均值为 $\widehat{\E}[Y\given A=90]=24.55+90\times2.14=216.9$。由于普通最小二乘估计使用全部数据点来寻找最佳直线，$A=a$ 组中 $Y$ 的均值 $\E[Y\given A=a]$ 是借用治疗取值不等于 $a$ 的个体的信息来估计的。

\begin{figure}[htbp]
\centering
\begin{tikzpicture}[x=0.062cm,y=0.0152cm,causal]
  \draw[->] (0,0) -- (0,272);
  \draw[->] (0,0) -- (104,0);
  \foreach \y in {50,100,150,200,250}{%
    \draw (-1.5,\y) -- (0,\y);
    \node[left=1pt,font=\footnotesize,inner sep=1.5pt] at (0,\y) {\y};}
  \node[left=1pt,font=\footnotesize,inner sep=1.5pt] at (0,0) {0};
  \foreach \x in {0,10,20,30,40,50,60,70,80,90,100}{%
    \draw (\x,0) -- (\x,-8);
    \node[below=1pt,font=\scriptsize,inner sep=0.8pt] at (\x,-8) {\x};}
  \node[below=13pt,font=\small] at (50,-8) {$A$};
  \foreach \xa/\ya in {2/22,11/55,13/10,18/33,24/102,30/85,38/62,44/157,%
                       44/188,53/120,59/228,66/112,70/215,79/200,84/140,97/218}{%
    \fill (\xa,\ya) circle[radius=1.7pt];}
  \draw[thick] (0,24.55) -- (100,{24.55+2.14*100});
\end{tikzpicture}
\caption{图 11.3 的数据与两参数线性模型的最小二乘直线。}
\label{fig:11.4}
\end{figure}

那么，什么是模型？模型由对数据联合分布的一种先验限制来定义。我们的线性条件均值模型说条件均值函数 $\E[Y\given A]$ 是一条直线，这就限制了它的形状。例如，该模型把 $A=90$ 时 $Y$ 的均值限制在 $A=80$ 时与 $A=100$ 时 $Y$ 的均值之间。这一限制由参数 $\theta_0$ 与 $\theta_1$ 编码。参数条件均值模型正是通过其先验限制来补充信息，以弥补数据中信息之不足。

参数估计量——基于参数条件均值模型的估计量——使我们能估计本来无法估计的量，例如当研究总体中不存在治疗水平 $A=90$ 的个体时，目标总体中这类个体的 $Y$ 均值。但天下没有免费的午餐。使用参数模型时，推断只有在模型所编码的限制正确、即模型被正确设定时才正确。于是，基于模型的因果推断——本书余下的大部分篇幅都致力于此——依赖于“（近似）无模型错误设定”这一条件。由于参数模型极少（如果有的话）被完美设定，一定程度的模型错误设定几乎总是可以预期的。这至少可以部分地通过使用非参数估计量来纠正，下一节即介绍。

\section{条件均值的非参数估计量}

回到图 11.1 的数据。治疗 $A$ 是二值的，我们希望一致地估计受治疗者的 $Y$ 均值 $\E[Y\given A=1]$ 与未受治疗者的 $\E[Y\given A=0]$。假设我们太迷恋模型，决定用模型来估计这两个量。我们再次提出线性模型
\[
\E[Y\given A]=\theta_0+\theta_1 A
\]
其中 $\E[Y\given A=0]=\theta_0+0\times\theta_1=\theta_0$，$\E[Y\given A=1]=\theta_0+1\times\theta_1=\theta_0+\theta_1$。我们用最小二乘法得到参数 $\theta_0$ 与 $\theta_1$ 的估计：$\hat\theta_0=67.5$、$\hat\theta_1=78.75$。于是我们估计 $\widehat{\E}[Y\given A=0]=67.5$、$\widehat{\E}[Y\given A=1]=146.25$。注意，基于模型的估计与 11.1 节算出的样本均值完全相同。这不是巧合，而是必然的结果。

\mnote{代码：程序 11.2。}

再来看治疗 $A$ 为二值时的模型 $\E[Y\given A]=\theta_0+\theta_1 A$。若把模型改写为 $\E[Y\given A=1]=\E[Y\given A=0]+\theta_1$，就会看到它只是说受治疗者的均值 $\E[Y\given A=1]$ 等于未受治疗者的均值 $\E[Y\given A=0]$ 加上一个量 $\theta_1$，而 $\theta_1$ 可正、可负、可为零。可是这句话当然永远成立！模型对 $\E[Y\given A=1]$ 与 $\E[Y\given A=0]$ 的取值没有施加任何限制。因此，$A$ 为二值时，$\E[Y\given A=a]=\theta_0+\theta_1 A$ 不是模型，因为它和样本均值一样，让数据自己说话。\mnote{本书把“模型”定义为对可能自然状态的先验数学限制（Robins, Greenland 1986）。第一部分题为“无模型的因果推断”，因为它只描述了饱和模型。}不对数据的分布施加限制的“模型”称为\emph{饱和模型}。饱和模型形式上像模型，尽管并不符合我们对模型的定义，因此通常也被称作模型。

一般地，当条件均值模型中的参数个数等于总体中未知条件均值的个数时，模型就是饱和的。例如，线性模型 $\E[Y\given A]=\theta_0+\theta_1 A$ 有两个参数，而 $A$ 为二值时总体中有两个未知条件均值：$\E[Y\given A=1]$ 与 $\E[Y\given A=0]$。既然模型不限制这两个参数的取值，也就不限制这两个均值的取值。\mnote{饱和模型等号两边的未知量个数相同。}作为对比，考虑图 11.3 中 $A$ 可取 0 到 100 的数据。线性模型 $\E[Y\given A]=\theta_0+\theta_1 A$ 只有两个参数，却要估计 101 个量，即 $\E[Y\given A=0],\E[Y\given A=1],\dots,\E[Y\given A=100]$。要用这两个参数无偏地估计这 101 个量，唯一的希望是碰巧这 101 个均值 $\E[Y\given A=a]$ 都落在一条直线上。当一个模型只有少数几个参数、却用来估计许多总体量时，我们说这个模型是\emph{简约的}（parsimonious）。

\begin{finepoint}[Fisher 一致性]
正文中对非参数估计量的定义，与统计学中所谓的\emph{Fisher 一致估计量}（Fisher 1922）吻合。也就是说，这种估计量若用整个总体而非样本计算，就给出总体参数的真值。按定义，Fisher 一致估计量不含任何模型限制；但如正文所述，许多总体量的 Fisher 一致估计可能并不存在。从技术上说，Fisher 一致估计量（当它们存在时）是饱和模型下总体量的非参数最大似然估计量。

在统计学文献中，“非参数估计量”一词有时指并非 Fisher 一致、但只施加极弱限制的估计量，例如核回归模型。细节见技术注记 11.1。
\end{finepoint}

我们在这里把条件均值函数的非参数估计量定义为：不对条件均值函数施加任何先验限制、直接由数据产生估计的估计量（更严格的定义见精要 11.1）。对二值治疗，总体均值 $\E[Y\given A=a]$ 的非参数估计量之一是其经验版本，即样本均值，或等价地，本节所述的饱和模型。当 $A$ 为有 100 个水平的离散变量、而样本中无人取 $A=90$ 时，$\E[Y\given A=90]$ 的非参数估计量不存在。本书第一部分介绍的所有因果推断方法——标准化、逆概率加权、分层、匹配——都基于饱和模型下总体量的非参数估计量，因为它们不对效应估计的取值施加任何先验限制。相比之下，本书第二部分介绍的大多数因果推断方法，所依赖的估计量是数据分布某一部分的参数估计量。当数据无法自己说话时（这很常见），参数估计就是借用信息的一种途径。

\mnote{可识别性假设是我们即便拥有无限多的数据也必须作出的假设，以便计算参数；建模假设则是由于我们没有无限多的数据、为估计参数而必须额外作出的假设。形式上，可识别性假设使参数成为观测数据联合分布的唯一函数。}

\section{平滑}

再次考虑图 11.3 的数据与线性模型 $\E[Y\given A]=\theta_0+\theta_1 A$。参数 $\theta_1$ 是治疗剂量 $A$ 每增加一个单位时结局均值的差。由于 $\theta_1$ 只是一个数，模型就设定：治疗 $A$ 每增加一个单位所致的结局 $Y$ 均值之差，在 $A$ 的整个取值范围内必须恒定，即模型要求条件均值结局随治疗剂量 $A$ 沿一条直线变化。图 11.4 给出了拟合最好的直线。

但可以设想这样的情形：在低剂量处，一个单位的变化带来的结局均值之差较大；在高剂量处则较小。若治疗剂量达到某一水平后，更高剂量的效应越来越小，就会出现这种情形。在这种情形下，模型 $\E[Y\given A]=\theta_0+\theta_1 A$ 是错的。不过，线性模型可以做得更灵活。

\mnote{注意：“线性”一词常有两种不同含义。当模型被表达为参数与变量函数的线性组合时，它就是线性的，即便后者是协变量的非线性函数（例如更高次幂或对数）。}

例如，设我们对图 11.3 的数据拟合模型 $\E[Y\given A]=\theta_0+\theta_1 A+\theta_2 A^2$，其中 $A^2=A\times A$ 即 $A$ 的平方。这仍称为线性模型，因为条件均值被表达为线性组合，即每个协变量（$A$ 与 $A^2$）与其相应系数（参数 $\theta_1$ 与 $\theta_2$）之积的和，再加上截距（$\theta_0$）。然而，只要 $\theta_2$ 不为零，$(\theta_0,\theta_1,\theta_2)$ 定义的就不再是直线，而是一条曲线——抛物线。我们把 $\theta_1$ 称为线性项 $A$ 的参数，把 $\theta_2$ 称为二次项 $A^2$ 的参数。

三参数线性模型 $\E[Y\given A]=\theta_0+\theta_1 A+\theta_2 A^2$ 下的曲线，可对图 11.3 的数据施用普通最小二乘估计得到，估计曲线见图 11.5。参数估计为 $\hat\theta_0=-7.41$、$\hat\theta_1=4.11$、$\hat\theta_2=-0.02$。治疗水平 $A=90$ 者中 $Y$ 的预测均值由下式给出：$\widehat{\E}[Y\given A=90]=\hat\theta_0+90\hat\theta_1+90\times90\hat\theta_2=197.1$。

\mnote{代码：程序 11.3。在同方差假设下，$\E[Y\given A=90]$ 的 Wald 95\% 置信区间为 $(142.8,\,251.5)$。}

我们可以不断加入参数：三次项（$\theta_3A^3$）、四次项（$\theta_4A^4$）……直到 15 次项（$\theta_{15}A^{15}$）。此时模型中的参数个数等于数据点（个体）数。曲线的形状会随参数个数增加而变化。一般地，模型中的参数越多，出现的拐点就越多。

也就是说，参数越多，曲线通常越“扭”，越不平滑。只有 2 个参数的线性模型——一条直线——是最平滑的模型。参数个数与数据点一样多的线性模型是最不平滑的模型，因为它的可能拐点数与数据点数一样多。事实上，这样的模型对数据做了插值，即样本中每个数据点都落在估计的条件均值函数上。

建模常常可以看作把噪声数据变换成某种平滑程度之曲线的过程。这种平滑之所以发生，是因为模型借用许多数据点的信息来预测协变量取某一特定组合时的结局值。平滑来自这样一点：$\E[Y\given A=a]$ 是借用 $A$ 不等于 $a$ 的个体的信息来估计的。所有参数估计量都在某种程度上纳入了平滑。

平滑的程度取决于在个体之间借用了多少信息。两参数模型 $\E[Y\given A]=\theta_0+\theta_1 A$ 借用研究总体中全部个体的信息来找出最小二乘直线，以估计 $\E[Y\given A=90]$。参数个数与个体数一样多的模型，在估计数据中出现的那些 $A$ 值处的 $\E[Y\given A]$ 时不借用任何信息，但对数据中未出现的 $A$ 值，它（通过插值）借用信息。\mnote{图 11.5}

中等程度的平滑可以通过使用中等数目的参数、或更一般地通过限制参与估计的个体数来实现。例如，为估计 $\E[Y\given A=90]$，我们可以决定只用治疗剂量在 80 与 100 之间的个体拟合两参数模型 $\E[Y\given A]=\theta_0+\theta_1 A$。也就是说，我们只借用 $A=90$ 附近一个 10 单位窗口内个体的信息；窗口越宽，平滑的程度就越高。

\begin{figure}[htbp]
\centering
\begin{tikzpicture}[x=0.062cm,y=0.0152cm,causal]
  \draw[->] (0,0) -- (0,272);
  \draw[->] (0,0) -- (104,0);
  \foreach \y in {50,100,150,200,250}{%
    \draw (-1.5,\y) -- (0,\y);
    \node[left=1pt,font=\footnotesize,inner sep=1.5pt] at (0,\y) {\y};}
  \node[left=1pt,font=\footnotesize,inner sep=1.5pt] at (0,0) {0};
  \foreach \x in {0,10,20,30,40,50,60,70,80,90,100}{%
    \draw (\x,0) -- (\x,-8);
    \node[below=1pt,font=\scriptsize,inner sep=0.8pt] at (\x,-8) {\x};}
  \node[below=13pt,font=\small] at (50,-8) {$A$};
  \foreach \xa/\ya in {2/22,11/55,13/10,18/33,24/102,30/85,38/62,44/157,%
                       44/188,53/120,59/228,66/112,70/215,79/200,84/140,97/218}{%
    \fill (\xa,\ya) circle[radius=1.7pt];}
  \draw[thick] plot[domain=0:100,samples=61] (\x,{-7.41+4.11*\x-0.02*\x*\x});
\end{tikzpicture}
\caption{图 11.3 的数据与三参数线性模型（含二次项）的最小二乘曲线。}
\label{fig:11.5}
\end{figure}

\mnote{我们用连续结局的模型作为例子。同样的推理适用于二值结局的模型，例如 logistic 模型（见技术注记 11.1）。}

在上面的简单例子中，所有模型都只含一个协变量（或对 $A$ 用一个参数，或对 $A$ 与 $A^2$ 用两个参数），因而曲线可以画在书页的二维平面上。在实际应用中，模型常包含许多不同的协变量，曲线实际上是超维曲面。无论问题的维数如何，平滑的概念不变：模型中的参数越少，预测（反应）曲面越平滑。

\section{偏倚—方差的权衡}

前几节我们用图 11.3 的 16 个个体，估计目标总体中接受治疗剂量 $A=90$ 者的结局均值 $\E[Y\given A=90]$。由于研究总体中无人接受 $A=90$，我们无法让数据自己说话，于是把数据与一个线性模型结合起来。估计 $\widehat{\E}[Y\given A=90]$ 随模型而变：在两参数模型 $\E[Y\given A]=\theta_0+\theta_1 A$ 下，估计值为 216.9（95\% 置信区间：172.1，261.6）；在三参数模型 $\E[Y\given A]=\theta_0+\theta_1 A+\theta_2 A^2$ 下，估计值为 197.1（95\% 置信区间：142.8，251.5）。我们用了两个不同的参数模型，得到两个不同的估计。哪一个更好？216.9 与 197.1 哪一个更接近目标总体中的均值？

若关系确实弯曲，则两参数模型的估计有偏，因为该模型假定了一条直线。另一方面，若关系确实是一条直线，则两个模型的估计都有效。这是因为无论关系是直线（此时 $\theta_2=0$）还是抛物线（此时 $\theta_2\neq0$），三参数模型 $\E[Y\given A]=\theta_0+\theta_1 A+\theta_2 A^2$ 都被正确设定。一个稳妥的策略是用三参数模型而非两参数模型。由于三参数模型在直线与抛物线两种情形下都被正确设定，它更不容易有偏。一般地，模型中的参数越多，模型施加的限制越少；模型越不平滑，对模型错误设定所致偏倚的防护就越强。

不平滑的模型可能给出偏倚更小的估计，但也会带来更大的方差，即估计周围更宽的 95\% 置信区间。例如，用三参数模型时 $\widehat{\E}[Y\given A=90]$ 周围的 95\% 置信区间比用两参数模型时宽得多。然而，当基于两参数模型的估计 $\widehat{\E}[Y\given A=90]$ 有偏时，标准的（名义）95\% 置信区间并未被校准，即它覆盖真参数 $\E[Y\given A=90]$ 的时间比例不是 95\%。

这种偏倚—方差的权衡是许多数据分析的核心。使用模型的研究者必须决定：为防偏倚（比如说，给模型增加参数）所付的方差代价是否值得。\mnote{精要 11.2 讨论模型维数对频率派区间与贝叶斯区间之关系的影响。}虽然有一些正规程序可辅助这类决策，实践中许多研究者仍依据传统、参数的可解释性、软件是否可得等标准来选择模型。本书通常假定我们的参数模型被正确设定。这是一个不现实的假设，但它使我们能专注于因果分析特有的问题。毕竟，模型错误设定是任何数据分析都可能出现的问题，无论估计值是否被赋予因果解释。在实践中，细心的研究者总会质疑其模型的有效性，并在与现有专家知识相容的不同模型设定下做替代分析。他们的目标是评估估计对模型设定的敏感性。

下面我们就可以讨论模型在因果推断中的用法了。

\begin{finepoint}[模型维数与频率派区间、贝叶斯区间的关系]
在频率派统计推断中，概率被定义为频率。在贝叶斯推断中，概率被定义为相信程度——一个与“概率即频率”很不相同的概念（de Finetti 1972）。第 10 章介绍了频率派统计推断所用的置信区间。贝叶斯统计推断则使用可信区间，其解释更为自然：一个贝叶斯 95\% 可信区间意味着，给定观测数据，“估计目标落在该区间内的概率为 95\%”。然而，部分由于需要设定研究者的相信程度，贝叶斯推断不如频率派推断常用。

有趣的是，在样本量很大的简单低维参数模型中，95\% 贝叶斯可信区间同时也是 95\% 频率派置信区间；而在高维模型或非参数模型中，贝叶斯 95\% 可信区间可能不是 95\% 置信区间，因为它覆盖估计目标的时间比例可能远低于 95\%。这些结果的根本原因在于，贝叶斯推断要求为所有未知参数指定先验分布。在低维参数模型中，数据中的信息会淹没合理先验所含的信息，因此推断对所选的具体先验分布相对不敏感。但在高维模型中就不再如此。于是，若生成数据的真实参数值在所选先验下不太可能出现，贝叶斯可信区间的中心就会被拉离真实参数，而偏向先验下概率最大的参数值。
\end{finepoint}

\begin{techpoint}[常用模型的分类]
正文描述的是形如
\[
\E[Y\given X]=\theta X\equiv\sum_{i=0}^{p}\theta_i X_i
\]
的线性条件均值模型，其中 $X$ 是协变量向量 $X_0,X_1,\dots,X_p$，且对所有 $n$ 个个体都有 $X_0=1$。这类模型是更大一类条件均值模型（McCullagh and Nelder, 1989; McCulloch, Searle, and Neuhaus, 2008）的子集，后者有两个组成部分：线性泛函形式或预测子 $\sum_{i=0}^{p}\theta_iX_i$，以及连结函数 $g\{\cdot\}$，使得 $g\{\E[Y\given X]\}=\sum_{i=0}^{p}\theta_iX_i$。

正文所述的线性条件均值模型使用恒等连结函数。取值严格为正的结局（例如计数、发病率的分子）的条件均值模型常用对数连结函数，以确保所有预测值都大于零，即 $\log\{\E[Y\given X]\}=\sum_{i=0}^{p}\theta_iX_i$，从而 $\E[Y\given X]=\exp\bigl(\sum_{i=0}^{p}\theta_iX_i\bigr)$。

二值结局（即只取 0 与 1 的结局）的条件均值模型常用 logit 连结，即 $\log\bigl\{\E[Y\given X]/(1-\E[Y\given X])\bigr\}=\sum_{i=0}^{p}\theta_iX_i$，从而 $\E[Y\given X]=\mathrm{expit}\bigl(\sum_{i=0}^{p}\theta_iX_i\bigr)$。这一连结确保所有预测值都大于 0 且小于 1。使用 logit 函数的条件均值模型称为 logistic 回归模型，本书广泛使用。对这些连结（称为典则连结），我们可以用最大似然估计 $\theta$：恒等连结下假定正态工作模型，对数连结下假定 Poisson 工作模型，logit 连结下假定 logistic 回归模型。只要 $\E[Y\given X]$ 的条件均值模型正确，这些估计对 $\theta$ 就是一致的。广义估计方程（GEE）模型常用于处理重复测量，是条件均值模型的又一个例子（Liang and Zeger, 1986）。

条件均值模型只规定 $\E[Y\given X]$ 的参数形式，并不另外限制 $Y\given X$ 的分布或 $X$ 的边缘分布。因此，当 $X$ 或 $Y$ 连续时，参数条件均值模型是数据 $(X,Y)$ 联合分布的\emph{半参数模型}，因为分布的一部分被参数化建模，另一部分则不受限制。之所以是半参数的，是因为联合分布中所有不受限制的成分不能用有限个参数表示。

条件均值模型本身也可以通过放松“$\E[Y\given X]$ 取参数形式”这一假定来推广。例如，核回归模型不对 $\E[Y\given X]$ 施加具体的函数形式，而是用
\[
\E[Y\given X=x]=\frac{\sum_{i=1}^{n}w_h(x-X_i)\,Y_i}{\sum_{i=1}^{n}w_h(x-X_i)}
\]
对任意 $x$ 作估计，其中 $w_h(z)$ 是一个取正值的函数，称为核函数，它在 $z=0$ 处达到最大值，并随 $|z|$ 变大而减小到 0，减小的速率取决于下标 $h$ 这一参数。另一个例子是广义可加模型（GAM），它把条件均值模型中的线性组合 $\sum_{i=0}^{p}\theta_iX_i$ 换成平滑函数之和 $\sum_{i=0}^{p}f_i(X_i)$。该模型可用回拟算法估计，其中 $f_i(\cdot)$ 在第 $k$ 次迭代时用（例如）核回归估计（Hastie and Tibshirani 1990）。

正文讨论的是用参数模型作平滑，这类模型为 $\E[Y\given X=x]$ 规定了先验的函数形式，例如抛物线。在估计 $\E[Y\given X=x]$ 时，模型可能借用离 $x$ 很远的 $X$ 值的信息。相比之下，核回归模型不规定先验函数形式，在估计 $\E[Y\given X=x]$ 时只借用 $x$ 附近 $X$ 值的信息。核回归模型是“非参数”回归模型的一例。这里“非参数”一词的用法与我们此前的用法不同：我们此前对 $\E[Y\given X=x]$ 的非参数估计量只使用 $X$ 恰好等于 $x$ 的个体，连近邻的信息都不借用；而这里的“非参数”估计量使用 $X$ 接近 $x$ 的个体，多近由称为带宽 $h$ 的平滑参数控制。
\end{techpoint}
